\documentclass[11pt]{article}

\usepackage[T1]{fontenc}
\usepackage{lmodern}
\usepackage{amsmath,amssymb,amsthm,mathtools}
\usepackage[margin=1in]{geometry}
\usepackage{microtype}
\usepackage{xcolor}
\usepackage{enumitem}
\usepackage{fancyhdr}
\usepackage[colorlinks=true,linkcolor=blue!45!black,citecolor=blue!45!black,urlcolor=blue!55!black]{hyperref}

\allowdisplaybreaks
\setlength{\parindent}{1.25em}
\setlength{\parskip}{0.15em}
\setlist{nosep}

\newtheorem{theorem}{Theorem}[section]
\newtheorem{lemma}[theorem]{Lemma}
\newtheorem{proposition}[theorem]{Proposition}
\newtheorem{corollary}[theorem]{Corollary}
\theoremstyle{definition}
\newtheorem{definition}[theorem]{Definition}
\theoremstyle{remark}
\newtheorem{remark}[theorem]{Remark}

\newcommand{\Z}{\mathbb{Z}}
\newcommand{\N}{\mathbb{N}}
\newcommand{\one}{\mathbf{1}}
\newcommand{\chips}[1]{\lvert #1\rvert}

\pagestyle{fancy}
\fancyhf{}
\fancyhead[L]{The middle stair in parallel chip-firing}
\fancyhead[R]{Draft for specialist review}
\fancyfoot[C]{\thepage}
\renewcommand{\headrulewidth}{0.4pt}
\fancypagestyle{plain}{%
  \fancyhf{}%
  \fancyhead[L]{The middle stair in parallel chip-firing}%
  \fancyhead[R]{Draft for specialist review}%
  \fancyfoot[C]{\thepage}%
  \renewcommand{\headrulewidth}{0.4pt}%
}

\title{The Middle Stair in Parallel Chip-Firing}
\author{%
  OpenAI Codex (GPT-5.6 Sol)\thanks{The documented model slug is
  \texttt{gpt-5.6-sol}. The exact runtime build identifier was not exposed to
  this session.}\\[-1pt]
  \small AI research system
  \and
  Nathan Lannan\\[-1pt]
  \small \href{mailto:nathan@nathanlannan.com}{\texttt{nathan@nathanlannan.com}}\\[-1pt]
  \small No institutional affiliation
}
\date{13 July 2026}

\begin{document}
\maketitle
\thispagestyle{empty}

\begin{abstract}
Let $G=(V,E)$ be a finite connected simple graph. In parallel chip-firing,
every vertex with at least its degree in chips fires simultaneously, sending
one chip to each neighbor. We prove that every nonnegative integral
configuration $\sigma$ satisfying
\[
   2|E|-|V|<\chips{\sigma}<2|E|
\]
has least eventual state period two. This is the statement posed as
Conjecture~2 by Ji, Li, and Wang. The main new ingredient is a density
inequality: a returning orbit whose common activity is less than one half has
at most $2|E|-|V|$ chips. Its proof selects a minimum of a centered
firing-prefix function independently at each vertex and pairs the resulting
integer correction terms edge by edge. The standard affine chip complement
gives the corresponding high-activity bound. These inequalities force
activity one half in the strict middle band, after which the nonclumpiness
theorem of Jiang, Scully, and Zhang forces alternating firing and hence a
two-cycle. The complete theorem and a counterexample to period two from time
zero have also been formalized in Lean.
\end{abstract}

\section{Introduction}

Parallel chip-firing is a deterministic synchronous dynamical system on
integer configurations of a graph. Its eventual behavior has been studied
since Bitar and Goles~\cite{BitarGoles1992}; see also
Jiang~\cite{Jiang2010}. On complete graphs, Levine connected the activity
function with a devil's staircase and a rotation number
\cite{Levine2011}. Recent work of Ji, Li, and Wang isolates the middle stair
for a general finite connected graph~\cite{JiLiWang2024}. They pose the
following assertion as their Conjecture~2.

\begin{theorem}[Middle-stair theorem]\label{thm:main}
Let $G=(V,E)$ be a finite connected simple graph and let
$\sigma:V\to\Z_{\ge 0}$ be a chip configuration. If
\[
   2|E|-|V|<\sum_{v\in V}\sigma(v)<2|E|,
\]
then the least eventual state period of $\sigma$ is exactly two.
\end{theorem}

The theorem was previously known in several graph classes; Ji, Li, and Wang
summarize the tree, cycle, complete-graph, and balanced complete-bipartite
cases~\cite{JiLiWang2024}. Our contribution is a graph-independent density
inequality for low-activity returning orbits (Lemma~\ref{lem:low-density}).
The complementary inequality follows from the established affine complement
for parallel chip-firing~\cite{Jiang2010}. These inequalities force activity
$1/2$. The remaining step uses the theorem of Jiang, Scully, and Zhang that a
periodic firing word cannot contain both $00$ and $11$
\cite[Theorem~6.2]{JiangScullyZhang2015}.

We make no priority claim beyond the literature review recorded with this
draft. Formal verification establishes correctness relative to the stated
definitions and dependencies, not originality. The low-activity density
inequality is therefore the main point requiring specialist comparison with
prior work.

\section{Dynamics, recurrence, and common activity}

Write $n=|V|$, $m=|E|$, and $d_v=\deg(v)$. A configuration is a map
$\sigma:V\to\Z_{\ge 0}$. At time $t$, define the firing indicator
\[
 f_t(v)=
 \begin{cases}
 1,&\sigma_t(v)\ge d_v,\\
 0,&\sigma_t(v)<d_v.
 \end{cases}
\]
All active vertices fire simultaneously. If $L=D-A$ is the graph Laplacian,
then one round is
\begin{equation}\label{eq:update}
   \sigma_{t+1}=\sigma_t-Lf_t.
\end{equation}
The total number of chips is conserved, and nonnegative configurations remain
nonnegative.

A positive integer $p$ is an \emph{eventual state period} of $\sigma_0$ if
there is $\mu\ge 0$ such that
\[
   \sigma_{\mu+p+t}=\sigma_{\mu+t}\qquad(t\ge 0).
\]
The least such $p$ is the least eventual state period. This is a period of the
state orbit, not merely of an aggregate activity statistic.

Since the coordinates are nonnegative and have fixed sum, every orbit is
eventually periodic. Restart at a state on a returning orbit and let $T>0$ be
any return length, so $\sigma_T=\sigma_0$. Put
\[
   u_v(t)=\sum_{s=0}^{t-1}f_s(v),\qquad 0\le t\le T.
\]
Telescoping~\eqref{eq:update} gives
\begin{equation}\label{eq:telescoping}
   \sigma_t=\sigma_0-Lu(t).
\end{equation}
At $t=T$ this yields $Lu(T)=0$. An integer-valued harmonic function on a
connected graph is constant, so there is an integer $q$, $0\le q\le T$, such
that
\begin{equation}\label{eq:common-count}
   u_v(T)=q\qquad\text{for every }v\in V.
\end{equation}
Thus every vertex has the same activity $q/T$. This common-count criterion is
standard in the subject; compare~\cite{BitarGoles1992,Jiang2010}.

We will use the following established firing-word theorem. Regard
$(f_0(v),\ldots,f_{T-1}(v))$ as a cyclic binary word.

\begin{theorem}[Jiang--Scully--Zhang nonclumpiness]\label{thm:nonclumpy}
On a returning parallel chip-firing orbit, the cyclic firing word of every
vertex does not contain both a cyclic occurrence of $00$ and a cyclic
occurrence of $11$.
\end{theorem}

This is the nonclumpiness theorem of
Jiang, Scully, and Zhang~\cite[Theorem~6.2]{JiangScullyZhang2015}. It enters
only after the density argument has forced $q/T=1/2$.

\section{A low-activity density inequality}

\begin{lemma}[Low-activity density bound]\label{lem:low-density}
Let $\sigma_0,\ldots,\sigma_T=\sigma_0$ be a returning orbit on a finite
connected simple graph. Let $q$ be the common firing count in
\eqref{eq:common-count}. If $2q<T$, then
\[
   \chips{\sigma_0}:=\sum_{v\in V}\sigma_0(v)\le 2m-n.
\]
\end{lemma}

\begin{proof}
For each vertex $v$, define the denominator-cleared centered prefix height
\[
   h_v(t)=T u_v(t)-qt,\qquad 0\le t\le T.
\]
Equation~\eqref{eq:common-count} gives $h_v(T)=h_v(0)=0$, so $h_v$ may be
viewed on the cyclic set of phases modulo $T$. Independently for every
vertex, select a phase at which $h_v$ is minimal. Let
$\tau_v\in\{0,\ldots,T-1\}$ be the round immediately preceding that phase.
When $\tau_v=T-1$, the endpoint $T$ represents phase zero.

Minimality against the preceding phase gives
\[
  0\ge h_v(\tau_v+1)-h_v(\tau_v)=T f_{\tau_v}(v)-q.
\]
Because $q<T$ and $f_{\tau_v}(v)\in\{0,1\}$, this forces
\begin{equation}\label{eq:selected-waits}
  f_{\tau_v}(v)=0.
\end{equation}

Fix an edge $\{v,w\}$. The independently selected minima satisfy
\[
  h_v(\tau_v+1)\le h_v(\tau_w),
  \qquad
  h_w(\tau_w+1)\le h_w(\tau_v).
\]
Using~\eqref{eq:selected-waits} and adding these inequalities gives
\begin{equation}\label{eq:edge-charge-bound}
   T K_{vw}\le 2q,
\end{equation}
where
\[
  K_{vw}=
  u_v(\tau_v)-u_w(\tau_v)
  +u_w(\tau_w)-u_v(\tau_w).
\]
The number $K_{vw}$ is an integer. Since $2q<T$, inequality
\eqref{eq:edge-charge-bound} implies
\begin{equation}\label{eq:edge-charge-sign}
   K_{vw}\le 0.
\end{equation}

Vertex $v$ waits at time $\tau_v$, so chip integrality gives
$\sigma_{\tau_v}(v)\le d_v-1$. Applying~\eqref{eq:telescoping} at the
vertex-dependent time $\tau_v$ yields
\[
  \sigma_0(v)
  \le d_v-1+(Lu(\tau_v))_v.
\]
Every one of these inequalities bounds a coordinate of the same initial
configuration, so they may be summed even though the selected times differ.
Regrouping the Laplacian terms by unoriented edges gives
\[
\begin{aligned}
  \chips{\sigma_0}
  &\le \sum_{v\in V}(d_v-1)
       +\sum_{\{v,w\}\in E}K_{vw}\\
  &\le 2m-n,
\end{aligned}
\]
where the last inequality uses~\eqref{eq:edge-charge-sign}. This proves the
claim.
\end{proof}

The bound is attained, for every connected graph with at least two vertices,
by the fixed configuration $\sigma(v)=d_v-1$. Thus its numerical threshold
cannot be improved without further hypotheses.

\section{Complementation and the high-activity bound}

\begin{lemma}[High-activity density bound]\label{lem:high-density}
Let $\sigma_0,\ldots,\sigma_T=\sigma_0$ be a returning orbit with common
firing count $q$. If $T/2<q<T$, then
\[
   \chips{\sigma_0}\ge 2m.
\]
\end{lemma}

\begin{proof}
Since $q<T$, every vertex waits at least once during the return interval. At
a time when $v$ waits, its chip count is at most $d_v-1$. The coordinatewise
upper bound
\begin{equation}\label{eq:periodic-box}
   \sigma_t(v)\le 2d_v-1
\end{equation}
then persists under one update: a waiting vertex receives at most $d_v$
chips, while a firing vertex loses $d_v$ chips and receives at most $d_v$.
Following the returning orbit once from a waiting phase for each $v$ proves
\eqref{eq:periodic-box} at every phase.

Define the affine complement
\[
   \sigma_t^c(v)=2d_v-1-\sigma_t(v).
\]
It is nonnegative by~\eqref{eq:periodic-box}. The threshold rule gives
$f_t^c=\one-f_t$, and $L\one=0$ shows that complementation commutes with the
update. Hence $\sigma^c$ is a returning orbit with common firing count
$T-q<T/2$. Lemma~\ref{lem:low-density} gives
$\chips{\sigma_0^c}\le 2m-n$. On the other hand,
\[
   \chips{\sigma_0^c}
   =\sum_{v\in V}(2d_v-1)-\chips{\sigma_0}
   =4m-n-\chips{\sigma_0}.
\]
Rearrangement gives $\chips{\sigma_0}\ge 2m$.
\end{proof}

The configuration $\sigma(v)=d_v$ is fixed because every vertex fires and
$L\one=0$; it attains the threshold $2m$. Together with the example following
Lemma~\ref{lem:low-density}, this explains why the inequalities in
Theorem~\ref{thm:main} must be strict.

\section{Proof of the middle-stair theorem}

\begin{proof}[Proof of Theorem~\ref{thm:main}]
Restart the eventual orbit at a returning state, and let $T>0$ be a return
length with common firing count $q$. The chip total remains in the strict
middle band.

The cases $q=0$ and $q=T$ are impossible. In the first case all vertices wait
throughout the return interval, so $\sigma(v)\le d_v-1$ and
$\chips{\sigma}\le 2m-n$. In the second case all vertices fire, so
$\sigma(v)\ge d_v$ and $\chips{\sigma}\ge 2m$.

If $2q<T$, Lemma~\ref{lem:low-density} contradicts the lower strict
inequality. If $T<2q$, then $q<T$ and Lemma~\ref{lem:high-density} contradicts
the upper strict inequality. Therefore
\begin{equation}\label{eq:half-activity}
   2q=T.
\end{equation}

Each cyclic firing word has $q$ ones and $T-q=q$ zeros. By
Theorem~\ref{thm:nonclumpy}, it has no cyclic $00$ or no cyclic $11$. A
balanced cyclic binary word with either property must alternate. Consequently
\[
   f_{t+1}(v)=1-f_t(v)
\]
for every vertex and every phase on the returning orbit. Applying
\eqref{eq:update} twice now gives
\[
\begin{aligned}
   \sigma_{t+2}
   &=\sigma_t-L(f_t+f_{t+1})\\
   &=\sigma_t-L\one
    =\sigma_t.
\end{aligned}
\]
Thus two is an eventual state period.

It remains to exclude eventual period one. At a fixed state, $Lf=0$; by
connectedness, the binary firing indicator $f$ is constant. If $f=0$, then
$\chips{\sigma}\le 2m-n$, and if $f=\one$, then
$\chips{\sigma}\ge 2m$. Conservation of the chip total contradicts the
strict band in either case. Hence one is not an eventual period, and the least
eventual state period is exactly two.
\end{proof}

\section{Eventuality and boundary cases}

The word ``eventual'' in Theorem~\ref{thm:main} cannot be removed. Let $G$ be
the path on three vertices, with the first coordinate at the center, and start
from
\[
   (0,0,2).
\]
Here $n=3$, $m=2$, and $1<2<4$. Direct application of the parallel update
gives
\[
 (0,0,2)\longmapsto(1,0,1)\longmapsto(2,0,0)
 \longmapsto(0,1,1)\longmapsto(2,0,0).
\]
Thus the initial state does not return after two rounds; the two-cycle begins
at time two. The displayed transitions and both period claims are checked in
the accompanying Lean development.

The endpoint configurations $d_v-1$ and $d_v$ are fixed and have totals
$2m-n$ and $2m$, respectively, when the connected graph has at least two
vertices. Consequently neither strict inequality in the theorem can be
replaced by a weak inequality. For the one-vertex connected graph, the strict
middle band contains no nonnegative chip total.

\section{Formal verification}

The accompanying Lean 4 development uses Lean 4.31.0 and mathlib v4.31.0.
Its terminal declaration is
\begin{quote}
\texttt{MiddleStair.middle\_stair\_least\_eventual\_period\_two}.
\end{quote}
The definitions distinguish returning cycles, eventual state periods, and
least positive eventual state periods. The development formalizes finite
recurrence, common firing counts, Lemmas~\ref{lem:low-density} and
\ref{lem:high-density}, half activity, alternation, the two-step return, and
the exclusion of eventual period one. It also formalizes the three-vertex
counterexample above.

For Theorem~\ref{thm:nonclumpy}, the development formalizes the sector and
mischief argument of Jiang, Scully, and Zhang. Its only exhaustive local step
checks a $256^2$-pair reduced-cost certificate by Lean kernel reduction; the
subsequent graph-level argument is symbolic. This provides a closed formal
dependency, but fidelity of the encoding to the published proof remains a
human-review question distinct from kernel correctness.

The project contains no \texttt{sorry}, \texttt{admit}, or declared project
axiom. An included \texttt{\#print axioms} audit reports only Lean and mathlib
logical foundations. Exact build instructions and a proof-to-code map are
included in the handoff repository.

\clearpage
\section*{Authorship and disclosure}

At Nathan Lannan's explicit direction, OpenAI Codex (GPT-5.6 Sol) is listed as
first author and Nathan Lannan as second author. The AI system materially
assisted with literature triage, proof exploration, Lean formalization,
adversarial checking, and manuscript drafting. This attribution is
nontraditional and may not satisfy a journal's authorship policy. Nathan
Lannan has no institutional affiliation and can be contacted at
\href{mailto:nathan@nathanlannan.com}{\texttt{nathan@nathanlannan.com}}. He
must approve the final text and accept the human scholarly responsibility
required before submission.

\bibliographystyle{plain}
\bibliography{references}

\end{document}
